
\PrintPage{III}
\chapter*{Introduction}

This volume aims to make $\ell$-adic cohomology easier for nonspecialists to
use. I hope that it will often enable them to avoid recourse to the dense
expos\'es of SGA~4 and SGA~5. It also contains some new results.

The first expos\'e, written by J.~F. Boutot, surveys SGA~4. It presents the
principal results---at a minimal level of generality that is often
insufficient for applications---and gives an idea of how they are proved. For
complete results or detailed proofs, SGA~4 remains indispensable.

The ``\hyperref[II]{Report on the trace formula}'' contains a complete proof of the trace formula for the
Frobenius endomorphism. The proof is the one Grothendieck gave in SGA~5,
pruned of every unnecessary detail. This Report should enable the user to
dispense with SGA~5, which may be regarded as a series of \var{digression}{digressions}{SGA45-CORR-0001}, some of
them very interesting. Its existence will allow SGA~5 to be published shortly
as it stands. It is supplemented by the expos\'e ``\hyperref[VI]{Applications of the trace formula to trigonometric sums},'' which
explains how the trace formula can be used to study trigonometric sums and
gives examples.

The intended readership of the other expos\'es is narrower, and their style
reflects that. The expos\'e ``\hyperref[III]{Functions $L$ modulo $\ell^n$ and modulo $p$}'' is a ``modular'' generalization
of the Report, based on the study of symmetric powers in SGA~4, Expos\'e~XVII,
\S5.5. The expos\'e ``\hyperref[IV]{The cohomology class associated with a cycle}'' defines this class in various contexts
and establishes compatibility between intersections and cup products. The
expos\'e ``\hyperref[V]{Duality}'' brings together some known results for which a
reference was lacking, together with some compatibility results. The expos\'e
``\hyperref[VII]{Finiteness theorems in $\ell$-adic cohomology}'' is new. In particular, for cohomology without supports it
gives finiteness theorems analogous to those known for cohomology with compact
supports.

For further details on the expos\'es, I refer to their respective
introductions.

\PrintPage{IV}
Finally, I thank J.~L. Verdier for allowing me to reproduce here his notes
``\hyperref[VIII]{Derived categories (State 0)}.'' I believe they remain very useful, and they had become
impossible to find.

In internal references within this volume, expos\'es are cited by an
abbreviated title, indicated in brackets in the table of contents.

\vspace{20pt}
Bures-sur-Yvette, September 20, 1976

\vspace{5pt}
Pierre Deligne

